% ---------------------------------------------------------------------------
% Section: Model  (input for theory_section.tex)
% ---------------------------------------------------------------------------
We describe the Rb gas with the ergodic quantum-kinetic theory of
Refs.~\cite{Zaremba1999,Bijlsma2000}, extended by the three processes that
define the experiment: forced evaporation in a trap of decreasing depth and
frequency, a loss that is uniform in energy and acts on both components
(collisions with the hot K atoms), and three-body recombination. Our choice of
this level of theory is a compromise between two considerations. First, the
quantity measured in the experiment is the condensed fraction of a cloud whose
atom number, temperature and trap frequency all change by factors of two to
four during the distillation, so the theory must evolve the full energy
distribution of the thermal cloud together with the condensate, and conserve
atom number and energy exactly in the absence of losses. The constant-temperature,
infinite-reservoir growth equations of Refs.~\cite{Gardiner1997,Gardiner1998}
cannot do this by construction. However, a full solution of the quantum
Boltzmann equation in phase space is not needed either: the elastic collision
rate ($\gamma_{\rm el}\simeq 10$ to $65$~Hz in the experiment) exceeds all other
rates in the problem by one to two orders of magnitude, so the distribution
function depends on position and momentum only through the single-particle
energy. This ergodic approximation reduces the kinetic equation to one
dimension and makes it possible to simulate the complete two-second protocol,
including its calibration, with modest computational resources.

\subsection{Kinetic equations}
The thermal atoms are described by the occupation number $g(\varepsilon,t)$ of
single-particle states of energy $\varepsilon$ in the Hartree-Fock potential
$U(\mathbf r,t)=V(\mathbf r)+2gn_c(\mathbf r,t)$. Here, $V=m\bar\omega^2r^2/2$ is
the external trap with geometric-mean frequency $\bar\omega(t)$,
$g=4\pi\hbar^2a/m$ is the coupling constant with $a=100.4\,a_0$
\cite{Egorov2013}, and $n_c$ is the condensate density, which we take in the
Thomas-Fermi form with chemical potential
$\mu_c=(\hbar\bar\omega/2)(15N_ca/l)^{2/5}$ and $l=\sqrt{\hbar/m\bar\omega}$. The
condensate is assumed to follow the number of condensed atoms $N_c(t)$
adiabatically~\cite{Bijlsma2000}; this is justified because the growth and
decay times found below ($\gtrsim 0.1$~s) are long compared with the trap
period. The thermal band starts at $\mu_c$, on the surface of the condensate,
where $U$ has its minimum~\cite{Bijlsma2000}; the mean-field repulsion of the
condensate therefore opens a gap between the condensate and the lowest thermal
states, which is the microscopic origin of the energy gap $\Delta$ invoked in
Ref.~\cite{MelladoMunoz2020}. We measure thermal energies from the band bottom,
$\bar\varepsilon=\varepsilon-\mu_c$, and write $\rho(\bar\varepsilon)$ for the
density of states, which is known analytically for the Thomas-Fermi profile
[Eqs.~(49)-(50) of Ref.~\cite{Bijlsma2000}] and reduces to
$\varepsilon^2/2(\hbar\bar\omega)^3$ without condensate.

The projected kinetic equation reads~\cite{Bijlsma2000}
\begin{equation}
\frac{\partial}{\partial t}(\rho g)+\frac{\partial}{\partial\bar\varepsilon}(\rho_w g)
= I_{22}+I_{12}-\gamma(t)\rho g-\Lambda_{\rm ev}\rho g ,
\label{eq:kinetic}
\end{equation}
with the collision integrals
\begin{align}
I_{22}(\bar\varepsilon_1)&=\frac{m^3g^2}{2\pi^3\hbar^7}\!\int\!d\bar\varepsilon_2d\bar\varepsilon_3d\bar\varepsilon_4\,
\rho(\bar\varepsilon_{\min})\,\delta(\bar\varepsilon_1+\bar\varepsilon_2-\bar\varepsilon_3-\bar\varepsilon_4)\nonumber\\
&\quad\times[(1+g_1)(1+g_2)g_3g_4-g_1g_2(1+g_3)(1+g_4)],
\label{eq:I22}\\
I_{12}(\bar\varepsilon_1)&=\frac{m^3g^2}{2\pi^3\hbar^7}N_c\!\int\!d\bar\varepsilon_3d\bar\varepsilon_4\,
\Phi(\bar\varepsilon_3,\bar\varepsilon_4)\,
[\delta(\bar\varepsilon_1-\bar\varepsilon_3-\bar\varepsilon_4)\nonumber\\
&\quad-\delta(\bar\varepsilon_1-\bar\varepsilon_3)-\delta(\bar\varepsilon_1-\bar\varepsilon_4)]\,
B(\bar\varepsilon_3,\bar\varepsilon_4),
\label{eq:I12}
\end{align}
where $g_i=g(\bar\varepsilon_i)$, $\bar\varepsilon_{\min}$ is the smallest of the
four energies, $B=(1+g_{3+4})g_3g_4-g_{3+4}(1+g_3)(1+g_4)$ with
$g_{3+4}=g(\bar\varepsilon_3+\bar\varepsilon_4)$, and $\Phi\in[0,1]$ is the fraction
of the condensate that can take part in a collision producing thermal atoms at
$\bar\varepsilon_3$ and $\bar\varepsilon_4$ [Eqs.~(57)-(60) of Ref.~\cite{Bijlsma2000}].
$I_{12}$ describes the collision of two thermal atoms that sends one of them
into the condensate and the inverse process; the number of condensed atoms
obeys $\dot N_c=-\int d\bar\varepsilon\,I_{12}-\gamma N_c-L_3\langle n_c^2\rangle N_c$.
Both directions of the exchange are contained in $B$: the condensate grows when
the thermal cloud is oversaturated ($g$ larger than the Bose distribution with
chemical potential $\mu_c$) and is depleted when it is undersaturated. Hence
the growth and loss rates $\bar W$ and $\bar K$ of the two-mode model of
Ref.~\cite{MelladoMunoz2020} are not parameters here but functionals of the
instantaneous distribution, and the gap enters through $\Phi$, which vanishes
quadratically when both thermal atoms have energies small compared with
$\mu_c$.

The term with the weighted density of states $\rho_w$ accounts for the shift
of the single-particle energies when the condensate mean field changes
[Eq.~(16) of Ref.~\cite{Bijlsma2000}]; we evaluate it as
$\rho_w=\dot m_c\,\partial_{m_c}\!\int_0^{\bar\varepsilon}\!\rho\,d\bar\varepsilon'$
with $m_c=\mu_c/\hbar\bar\omega$. Energies are measured in units of
$\hbar\bar\omega(t)$, in which an adiabatic change of the trap frequency leaves
the occupation numbers unchanged, so the evaporation ramp requires no further
term. Atom number is conserved exactly by Eqs.~(\ref{eq:kinetic})-(\ref{eq:I12})
without losses, and the total energy
$E=\int d\bar\varepsilon\,(\bar\varepsilon+\mu_c)\rho g+\tfrac57\mu_cN_c$ is conserved
up to the discretization error (0.3\% over 6~s in our benchmarks [SIM]).

\subsection{Evaporation, dissipation and losses}
Evaporation is modelled as in Ref.~\cite{Luiten1996}: states above the trap
depth $U_t(t)$ are empty and an atom scattered into one of them leaves the
trap, which appears in Eqs.~(\ref{eq:I22})-(\ref{eq:I12}) through final states
with $g=0$; atoms left above $U_t$ by the lowering of the trap are removed at
the trap frequency ($\Lambda_{\rm ev}$). For the crossed dipole trap we take
$U_t(t)=\eta_1k_BT_1[\bar\omega(t)/\bar\omega_1]^{p}$, with $T_1$ the measured
temperature at $t=1$~s and $\eta_1$, $p$ calibration parameters; $p=2$ for a
trap whose depth and frequency follow the beam power alone, and $p>2$ when
gravity reduces the effective depth. The frequency follows the ramp of
Ref.~\cite{MelladoMunoz2020}, $\bar\omega(t)=2\pi s_\omega(192-56\,t/{\rm s})$~Hz,
with $s_\omega$ a scale factor between the value calibrated from Eq.~(1) of
Ref.~\cite{MelladoMunoz2020} ($s_\omega=1$) and the value measured by parametric
heating ($s_\omega\simeq0.8$). The dissipation enters as a loss rate
$\gamma(t)=\gamma_K(t)+\gamma_{\rm bg}$ that is the same for every energy and
for the condensate: for the collision energies of the experiment the
momentum transfer in a Rb-K collision exceeds the trap depth in all but
$\sim10^{-3}$ of the events, so the collision removes the atom whatever its
state. We use $\gamma_K(t)=a\,t$ with $a=0.112$~Hz/s from Fig.~1(b) of
Ref.~\cite{MelladoMunoz2020}, and a background rate $\gamma_{\rm bg}$ fitted
on the data without K. Three-body recombination uses
$L_3=5.8\times10^{-30}$~cm$^6$/s for the condensate~\cite{Burt1997,Soding1999}
with $\langle n_c^2\rangle=\tfrac{8}{21}n_{c,{\rm peak}}^2$; for the thermal cloud
the corresponding rate is below $10^{-3}$~s$^{-1}$ at the densities of the
experiment and is neglected.

\subsection{Numerical method and observables}
Equation~(\ref{eq:kinetic}) is discretized on a uniform grid of $n=240$ to $400$
energies between the band bottom and $1.15\,U_t(t_0)$. On a uniform grid the
energy-conserving delta functions become index constraints, so both collision
integrals conserve number and energy to machine precision; the double sum of
Eq.~(\ref{eq:I22}) is evaluated in $O(n^2)$ operations by splitting it
according to which pair of states contains the lowest energy. Time integration
uses an explicit adaptive Runge-Kutta scheme. [SIM: grid convergence, 3\% in
$N_c$ between $n=200$ and 400.] The temperature of the thermal cloud is defined
from its mean energy per atom through the equilibrium relation of the same
model, which coincides with the Bose-distribution temperature whenever the
cloud is thermal; we also compute the temperature a Gaussian fit to the
momentum distribution would return, which is systematically lower (by 20\% at
saturation [SIM]), and use it to bound the systematic error of the comparison
with the Gaussian-fit temperatures of Ref.~\cite{MelladoMunoz2020}. The
chemical potential of the thermal cloud, $\tau=T/T_c$ and the equilibrium
fraction of Eq.~(1) are computed from the model output with the same
formulas and the same calibrated frequency as in Ref.~\cite{MelladoMunoz2020},
so that model and data are processed identically.
